\documentclass[12pt]{article}
\usepackage[utf8]{inputenc}
\usepackage[T1]{fontenc}
\usepackage[margin=1in]{geometry}
\usepackage{setspace}
\usepackage{xcolor}
\usepackage[colorlinks=true, urlcolor=blue!60!black, linkcolor=blue!60!black]{hyperref}
\setlength{\parindent}{2em}
\setlength{\parskip}{1.25em}
\renewcommand{\baselinestretch}{1.0}
\begin{document}
\begin{center}
{\LARGE \textit{Sherbert v. Verner}, 374 U.S. 398 (1963)\par}
\vspace{0.8em}
{\large Case Brief\par}
\end{center}
\vspace{1em}

\section{Memory Jogger}

In \href{https://www.courtlistener.com/opinion/106660/sherbert-v-verner/}{\textit{Sherbert v. Verner}}, the Supreme Court held that South Carolina violated the Free Exercise Clause by denying unemployment benefits to a Seventh-day Adventist who refused Saturday work on her Sabbath, establishing strict scrutiny (the ``compelling state interest'' test) for substantial burdens on religious exercise.

\section{Detailed Case Facts}

Appellant Adell H. Sherbert worked for approximately 35 years as a textile-mill spool tender for Spartan Mills (Beaumont Division) in Spartanburg, South Carolina (\href{https://law.justia.com/cases/south-carolina/supreme-court/1962/17915-1.html}{S.C. Supreme Court Opinion}). On August 5, 1957, she became a member of the Seventh-day Adventist Church, whose sincere and undisputed tenet forbids labor on the Sabbath from sundown Friday to sundown Saturday. At the time of her conversion, Spartan Mills permitted her to work a five-day week, Monday through Friday, with Saturday work optional (\href{https://supreme.justia.com/cases/federal/us/374/398/}{Justia Opinion}).

On June 5, 1959, Spartan Mills notified Sherbert that it was shifting all three shifts to a mandatory six-day workweek including Saturday. After Sherbert refused to report on six successive Saturdays for religious reasons, the mill discharged her on July 27, 1959. She sought employment at three other textile mills in the Spartanburg area and expressed willingness to work in another industry on a five-day schedule, but could not obtain non-Saturday work because nearly all local mills operated six days a week. (Of roughly 150 Seventh-day Adventists in the Spartanburg area, all but Sherbert and one other worker had secured non-Saturday employment.) (\href{https://supreme.justia.com/cases/federal/us/374/398/}{Justia Opinion}; \href{https://law.justia.com/cases/south-carolina/supreme-court/1962/17915-1.html}{S.C. Supreme Court Opinion}).

On July 29, 1959, Sherbert filed for unemployment compensation under the South Carolina Unemployment Compensation Act, S.C.\ Code \S\S~68-1 to 68-404. Under \S~68-113(3), a claimant must be ``able to work and is available for work,'' and under \S~68-114(3)(a), a claimant is disqualified if she fails ``without good cause'' to accept ``available suitable work'' when offered. Meanwhile, a separate South Carolina statute, S.C.\ Code \S~64-4, expressly protected Sunday worshippers in textile plants during declared national emergencies from being required to work on Sunday or penalized for conscientious objection to Sunday work (\href{https://tile.loc.gov/storage-services/service/ll/usrep/usrep374/usrep374398/usrep374398.pdf}{U.S. Reports PDF}).

\section{Procedural History}

\begin{itemize}
\item \textbf{South Carolina Employment Security Commission:} On September 4, 1959, a claims examiner found Sherbert unavailable for work and imposed a disqualification. The Appeal Tribunal affirmed on October 12, 1959, and the Full Commission (chaired by Charlie V. Verner) affirmed on December 18, 1959, ruling that her refusal of Saturday work constituted a failure without good cause to accept suitable work (\href{https://law.justia.com/cases/south-carolina/supreme-court/1962/17915-1.html}{S.C. Supreme Court Opinion}).
\item \textbf{Court of Common Pleas for Spartanburg County:} On January 5, 1960, Sherbert sought judicial review; on June 27, 1960, Presiding Judge J. Woodrow Lewis affirmed the Commission's disqualification order (\href{https://law.justia.com/cases/south-carolina/supreme-court/1962/17915-1.html}{S.C. Supreme Court Opinion}).
\item \textbf{Supreme Court of South Carolina (240 S.C. 286, 125 S.E.2d 737 (1962)):} Affirmed on May 17, 1962 (4--1, Justice Bussey dissenting on statutory grounds), holding that Sherbert's disqualification placed no restriction on her freedom of religion under the First and Fourteenth Amendments (\href{https://supreme.justia.com/cases/federal/us/374/398/}{Justia Opinion}).
\item \textbf{U.S. Supreme Court (374 U.S. 398 (1963)):} Noted probable jurisdiction (371 U.S. 938), heard oral argument on April 24, 1963, and reversed and remanded on June 17, 1963 (\href{https://www.courtlistener.com/opinion/106660/sherbert-v-verner/}{CourtListener}; \href{https://www.oyez.org/cases/1962/526}{Oyez}).
\end{itemize}

\section{Judicial Votes}

The Court decided the case \textbf{7--2} in favor of Sherbert (\href{https://supreme.justia.com/cases/federal/us/374/398/}{Justia Opinion}; \href{https://www.oyez.org/cases/1962/526}{Oyez}):
\begin{itemize}
\item \textbf{Majority Opinion (6 Justices):} Justice William J. Brennan, Jr., delivered the opinion of the Court, joined by Chief Justice Earl Warren and Justices Hugo L. Black, William O. Douglas, Tom C. Clark, and Arthur J. Goldberg.
\item \textbf{Concurring Opinion (1 Justice, also joining the majority):} Justice William O. Douglas filed a separate concurring opinion.
\item \textbf{Concurring in the Result (1 Justice):} Justice Potter Stewart filed an opinion concurring in the result only, declining to join the majority's reasoning.
\item \textbf{Dissenting Opinion (2 Justices):} Justice John Marshall Harlan II filed a dissenting opinion, joined by Justice Byron R. White.
\end{itemize}

\section{Holding}

\textbf{Judgment:} Reversed and remanded (\href{https://tile.loc.gov/storage-services/service/ll/usrep/usrep374/usrep374398/usrep374398.pdf}{U.S. Reports PDF}).

\textbf{Rule Announced:} Under the Free Exercise Clause of the First Amendment (applicable to the States via the Fourteenth Amendment), a State may not condition eligibility for public unemployment benefits in a manner that forces a worker to choose between following a cardinal tenet of her religious faith and forfeiting public benefits, unless the State proves that the substantial infringement is justified by a ``compelling state interest'' and that no less restrictive alternative form of regulation would achieve that interest without infringing First Amendment rights (\href{https://supreme.justia.com/cases/federal/us/374/398/}{Justia Opinion}).

\textbf{Holding vs.\ Dicta:} The Court's holding is that South Carolina's application of its unemployment-eligibility provisions to disqualify a Saturday Sabbatarian violated the Free Exercise Clause because it imposed a substantial indirect burden unsupported by a compelling state interest, and that extending benefits did not violate the Establishment Clause. The Court expressly declined to reach Sherbert's Fourteenth Amendment Equal Protection claim. Justice Brennan's statements that the Court was not declaring a constitutional right to unemployment benefits for persons whose religious convictions render them ``nonproductive member[s] of society,'' and his skeptical footnote and text questioning whether the state could constitutionally inquire into the truth or falsity of religious beliefs under \textit{United States v. Ballard}, 322 U.S. 78 (1944), are limiting dicta (\href{https://supreme.justia.com/cases/federal/us/374/398/}{Justia Opinion}).

\section{Analysis of Opinions}

\subsection{Majority Opinion (Justice Brennan)}

Justice Brennan framed the analysis in three steps. First, he concluded that disqualifying Sherbert imposed a substantial burden on her free exercise of religion. Although the law imposed no criminal penalty for Saturday worship, an indirect burden is constitutionally invalid when its effect is to impede religious observance. Forcing Sherbert to choose between obeying her Sabbath and forfeiting benefits operated identically to ``a fine imposed against appellant for her Saturday worship.'' Invoking the unconstitutional-conditions doctrine (\textit{Speiser v. Randall}, 357 U.S. 513 (1958)), Brennan rejected the argument that unemployment benefits are a mere ``privilege'' rather than a ``right.'' Moreover, South Carolina's statutory protection for Sunday worshippers under S.C.\ Code \S~64-4 compounded the constitutional infirmity by discriminating between faiths (\href{https://supreme.justia.com/cases/federal/us/374/398/}{Justia Opinion}).

Second, Justice Brennan applied strict scrutiny: a merely rational relationship to a colorable state interest would not suffice; only a ``compelling state interest'' addressing ``the gravest abuses, endangering paramount interests'' could justify the infringement. South Carolina's speculative fear of fraudulent claims diluting the unemployment fund or disrupting Saturday scheduling had neither been raised below nor supported by record evidence, and the State failed to show that no alternative regulatory means could prevent fraud without burdening religious liberty. Brennan distinguished \textit{Braunfeld v. Brown}, 366 U.S. 599 (1961) (which upheld Sunday closing laws against Orthodox Jewish merchants), reasoning that a uniform day of rest served a strong secular interest that would be rendered administratively unworkable by exemptions (\href{https://constitutioncenter.org/the-constitution/supreme-court-case-library/sherbert-v-verner}{National Constitution Center}).

Third, Justice Brennan rejected any Establishment Clause objection. Extending unemployment benefits to Sabbatarians on equal footing with Sunday worshippers merely reflected ``the governmental obligation of neutrality in the face of religious differences'' under \textit{Everson v. Board of Education}, 330 U.S. 1 (1947), rather than establishing the Seventh-day Adventist Church (\href{https://supreme.justia.com/cases/federal/us/374/398/}{Justia Opinion}).

\subsection{Concurring Opinion (Justice Douglas)}

Justice Douglas joined the Court's opinion and wrote separately to emphasize that the Free Exercise Clause fences off an individual's religious scruples from majoritarian pressure, illustrating the point with Muslim Friday prayers, Sikh ceremonial swords, Jehovah's Witness colporteurs, Quaker affirmations, and Buddhist dietary rules. Reiterating his dissent from \textit{Braunfeld v. Brown}, Douglas stressed that the harm was government interference with conscience, not measurable financial damage. He distinguished affirmative demands for state subsidies from nondiscriminatory welfare eligibility: because Sherbert would receive benefits as an unemployed worker rather than as an Adventist, no Establishment Clause problem arose (\href{https://supreme.justia.com/cases/federal/us/374/398/}{Justia Opinion}).

\subsection{Concurrence in the Result (Justice Stewart)}

Justice Stewart agreed that Sherbert was entitled to benefits, but refused to join the majority's opinion, identifying a ``double-barreled dilemma'' that the Court had ``paper[ed] over.'' First, Stewart argued that mandating unemployment benefits for a religious refusal to work Saturdays---while South Carolina denied benefits to a mother unavailable on Saturdays for lack of a babysitter---squarely collided with the Court's rigid, ``positively wooden'' Establishment Clause decisions in \textit{Engel v. Vitale}, 370 U.S. 421 (1962), and \textit{Abington School District v. Schempp}, 374 U.S. 203 (1963), which forbade any public financial support of religion. In Stewart's view, the Court should repudiate its mechanistic Establishment Clause rhetoric and recognize that the Free Exercise Clause requires ``an atmosphere of hospitality and accommodation.'' Second, Stewart argued that \textit{Sherbert} could not honestly be reconciled with \textit{Braunfeld v. Brown}, where a criminal statute threatened to destroy a merchant's entire business---a far heavier burden than losing up to 22 weeks of unemployment benefits. He therefore urged overruling \textit{Braunfeld} (\href{https://supreme.justia.com/cases/federal/us/374/398/}{Justia Opinion}).

\subsection{Dissenting Opinion (Justice Harlan, joined by Justice White)}

Justice Harlan argued that the South Carolina Unemployment Compensation Act was enacted during the Great Depression solely to relieve ``involuntary unemployment'' caused by industry's inability to provide jobs, and that South Carolina courts uniformly disqualified any worker unavailable for personal reasons---including childcare responsibilities. Because Sherbert was disqualified under the same neutral rule applied to secular personal reasons, the State had not discriminated against her religion. Harlan charged that the majority effectively overruled \textit{Braunfeld v. Brown} and constitutionally compelled the State to single out religiously motivated conduct for financial assistance while denying it to identically situated secular claimants. While Harlan believed a legislature \textit{could} permissibly accommodate Sabbatarians without violating the Establishment Clause, he insisted the Constitution did not \textit{compel} South Carolina to carve out a religious exemption from a neutral welfare statute (\href{https://supreme.justia.com/cases/federal/us/374/398/}{Justia Opinion}; \href{https://fedsoc.org/case/sherbert-v-verner}{Federalist Society}).

\section{Examples: Future Applications}

\begin{enumerate}
\item \textbf{Same-Side Application (Individualized ``Good Cause'' Unemployment Exemption):} A Jehovah's Witness quits his factory job after his employer transfers him from a sheet-metal division to a turret-manufacturing division making military tanks, which violates his sincere religious pacifism, and the state denies him unemployment benefits for quitting ``without good cause.'' \textit{Outcome:} Unconstitutional under \textit{Sherbert} (and its successor, \textit{Thomas v. Review Board}, 450 U.S. 707 (1981)). Conditioning unemployment benefits on abandoning a religious precept within an individualized ``good cause'' system imposes a substantial burden that fails strict scrutiny, an exception that survived even \textit{Employment Division v. Smith}.

\item \textbf{Same-Side Application (Discretionary Exemption Scheme Refusing Religious Accommodation):} A city contracts with private foster-care agencies under a nondiscrimination policy that empowers a city commissioner to grant individualized secular exemptions in her ``sole discretion,'' yet the city refuses to grant an exemption to a Catholic foster agency whose religious beliefs prevent it from certifying same-sex couples. \textit{Outcome:} Triggers strict scrutiny under \textit{Sherbert}'s individualized-exemption principle (as applied in \href{https://www.law.cornell.edu/supremecourt/text/19-123}{\textit{Fulton v. City of Philadelphia}, 593 U.S. 522 (2021)}). Because the system invites individualized assessments of the reasons for noncompliance, denying a religious hardship exemption is not generally applicable and must satisfy strict scrutiny.

\item \textbf{Opposite-Side Application (Neutral, Generally Applicable Criminal Prohibition):} Two state-licensed drug-rehabilitation counselors are fired and denied unemployment benefits after ingesting peyote during a religious ceremony in violation of a state criminal law that prohibits possession of controlled substances across the board without individualized exemptions. \textit{Outcome:} Constitutional under current Free Exercise doctrine (\href{https://fedsoc.org/case/employment-division-v-smith}{\textit{Employment Division v. Smith}, 494 U.S. 872 (1990)}). The Court refused to extend \textit{Sherbert}'s compelling-interest test to neutral, generally applicable criminal laws, reasoning that doing so would make every believer ``a law unto himself.''

\item \textbf{Opposite-Side Application (Generally Applicable Tax Law Without Individualized Exemptions):} An Amish employer refuses on religious grounds to withhold or pay federal Social Security taxes for his employees, arguing that his faith requires the community to care for its own elderly members. \textit{Outcome:} Constitutional (as in \textit{United States v. Lee}, 455 U.S. 252 (1982)). Even under the \textit{Sherbert} era's compelling-interest framework, the government's overriding interest in maintaining a comprehensive, uniform national tax system cannot function if myriad religious exemptions are constitutionally mandated.

\item \textbf{Fence-Sitter (Unemployment Benefits Denied to a Worker Who Refuses Mandatory Sunday Overtime for Secular Moral Conscience):} A nonreligious ethical humanist refuses mandatory Saturday shifts to care for a terminally ill parent out of deeply held secular moral duty, is fired, and challenges South Carolina's denial of unemployment benefits while Saturday Sabbatarians receive them. \textit{Outcome:} Genuine close call that highlights the exact tension Justices Stewart and Harlan flagged in \textit{Sherbert}. Under \textit{Sherbert} as written, Free Exercise protection attaches specifically to \textit{religious} beliefs, leaving secular conscientious objectors unprotected; yet favoring religious over comparably weighty secular moral obligations fuels ongoing scholarly and judicial debate over whether equal-regard or Establishment Clause principles forbid such asymmetry (\href{https://engagedscholarship.csuohio.edu/cgi/viewcontent.cgi?article=1116&context=clevstlrev}{Cleveland State Law Review (2009)}).
\end{enumerate}

\section{Critique}

\textbf{Originalist and Textualist Perspectives:} Originalist scholarship on \textit{Sherbert} is sharply divided over whether the Free Exercise Clause originally mandated religious exemptions from neutral, generally applicable laws. Critics such as Philip Hamburger and Justice Antonin Scalia (in \href{https://fedsoc.org/case/employment-division-v-smith}{\textit{Employment Division v. Smith}}) argue that Founding-era evidence does not support constitutionally compelled judicial exemptions from general laws, warning that \textit{Sherbert}'s strict scrutiny invites judges to weigh the centrality of religious beliefs and threatens the rule of law (\href{https://www.wakeforestlawreview.com/wp-content/uploads/2024/07/w05_Lingo.pdf}{Wake Forest Law Review (2024)}). Similarly, structural textualists argue that reading the First Amendment to compel affirmative financial accommodations for religion creates a collision with the Establishment Clause (\href{https://sclawreview.org/article/the-single-religion-clause-non-established-free-and-independent-exercise/}{South Carolina Law Review}). Conversely, originalists such as Michael McConnell contend that Founding-era state constitutions, colonial accommodations (e.g., for Quakers), and James Madison's writings support \textit{Sherbert}'s substantive-liberty view of religious exemptions, while Kurt Lash argues that the Fourteenth Amendment's Privileges or Immunities Clause supports robust religious accommodation against the States (\href{https://www.wakeforestlawreview.com/wp-content/uploads/2024/07/w05_Lingo.pdf}{Wake Forest Law Review (2024)}; \href{https://ir.lawnet.fordham.edu/cgi/viewcontent.cgi?article=3071&context=flr}{Fordham Law Review (1994)}).

\textbf{Progressive and Doctrinal Perspectives:} Doctrinal historians recognize \textit{Sherbert} as a watershed moment that helped birth modern strict scrutiny (\href{https://www.uclalawreview.org/wp-content/uploads/2019/09/33_54UCLALRev1267June2007.pdf}{UCLA Law Review (2007)}) and protected religious minorities from majoritarian indifference in the welfare state. After \textit{Smith} curtailed \textit{Sherbert} in 1990, Congress enacted the Religious Freedom Restoration Act of 1993 (RFRA) specifically to restore the \textit{Sherbert} compelling-interest test for federal action (\href{https://ir.lawnet.fordham.edu/cgi/viewcontent.cgi?article=3071&context=flr}{Fordham Law Review (1994)}). Yet many progressive and equality-oriented scholars---such as Christopher Eisgruber and Lawrence Sager---critique \textit{Sherbert}'s strict-scrutiny framework from a different angle: by treating religion as uniquely privileged over deep secular commitments, \textit{Sherbert}-style exemptions can undermine equal citizenship and, in modern litigation, permit religious claimants to seek exemptions from civil-rights and antidiscrimination laws that protect third parties (\href{https://engagedscholarship.csuohio.edu/cgi/viewcontent.cgi?article=1116&context=clevstlrev}{Cleveland State Law Review (2009)}).

\textbf{Underlying Competing Values:} \textit{Sherbert} crystallizes the enduring constitutional conflict between \textit{substantive religious liberty} (protecting minority faiths from indirect coercion by majoritarian norms) and \textit{formal neutrality and democratic governance} (ensuring that general welfare and regulatory statutes apply evenhandedly to religious and secular citizens alike without judicial ad hoc balancing).

\section{Key Quotations}

\begin{quote}
``The ruling forces her to choose between following the precepts of her religion and forfeiting benefits, on the one hand, and abandoning one of the precepts of her religion in order to accept work, on the other hand. Governmental imposition of such a choice puts the same kind of burden upon the free exercise of religion as would a fine imposed against appellant for her Saturday worship.''
--- Justice Brennan, Opinion of the Court, 374 U.S. at 404 (\href{https://www.courtlistener.com/opinion/106660/sherbert-v-verner/}{CourtListener}; \href{https://supreme.justia.com/cases/federal/us/374/398/}{Justia})
\end{quote}

\begin{quote}
``It is too late in the day to doubt that the liberties of religion and expression may be infringed by the denial of or placing of conditions upon a benefit or privilege.''
--- Justice Brennan, Opinion of the Court, 374 U.S. at 404 (\href{https://www.courtlistener.com/opinion/106660/sherbert-v-verner/}{CourtListener}; \href{https://supreme.justia.com/cases/federal/us/374/398/}{Justia})
\end{quote}

\begin{quote}
``It is basic that no showing merely of a rational relationship to some colorable state interest would suffice; in this highly sensitive constitutional area, `[o]nly the gravest abuses, endangering paramount interests, give occasion for permissible limitation,' \textit{Thomas v. Collins}, 323 U.S. 516, 530.''
--- Justice Brennan, Opinion of the Court, 374 U.S. at 406 (\href{https://www.courtlistener.com/opinion/106660/sherbert-v-verner/}{CourtListener}; \href{https://supreme.justia.com/cases/federal/us/374/398/}{Justia})
\end{quote}

\begin{quote}
``And I think that the guarantee of religious liberty embodied in the Free Exercise Clause affirmatively requires government to create an atmosphere of hospitality and accommodation to individual belief or disbelief.''
--- Justice Stewart, Concurring in the Result, 374 U.S. at 416 (\href{https://www.courtlistener.com/opinion/106660/sherbert-v-verner/}{CourtListener}; \href{https://supreme.justia.com/cases/federal/us/374/398/}{Justia})
\end{quote}

\begin{quote}
``The State, in other words, must single out for financial assistance those whose behavior is religiously motivated, even though it denies such assistance to others whose identical behavior (in this case, inability to work on Saturdays) is not religiously motivated.''
--- Justice Harlan, Dissenting, 374 U.S. at 422 (\href{https://www.courtlistener.com/opinion/106660/sherbert-v-verner/}{CourtListener}; \href{https://supreme.justia.com/cases/federal/us/374/398/}{Justia})
\end{quote}

\end{document}
